\chapter{System Architecture}
\label{ch:architecture}

The system is organized into three distinct layers, each with clearly defined responsibilities.

\section{Three-Layer Architecture}

\subsection{Layer 1 --- Physical Prototype}

The physical layer consists of an ESP32 microcontroller connected to energy meters (PZEM-004T for AC loads, INA226 for DC PV and battery measurements) and environmental sensors (DHT22, BH1750). All mains-voltage measurement is handled by isolated commercial meters; the ESP32 circuit operates exclusively at 3.3--5~V DC.

\subsection{Layer 2 --- Digital / Software Layer}

The software layer comprises:
\begin{itemize}[nosep]
    \item MQTT broker (Eclipse Mosquitto) for IoT communication.
    \item Python backend (FastAPI) for data ingestion, validation, and API exposure.
    \item Database (SQLite for development, PostgreSQL for production) for persistent storage.
    \item Data preprocessing pipeline for cleaning and feature engineering.
\end{itemize}

\subsection{Layer 3 --- Intelligence Layer}

The intelligence layer includes:
\begin{itemize}[nosep]
    \item ML forecasting models (consumption and PV production).
    \item Anomaly detection module.
    \item MILP optimization engine for battery dispatch and load scheduling.
    \item Streamlit dashboard for visualization and operator interaction.
\end{itemize}

\section{Data Flow}

The complete data flow from sensor to operator is:

\begin{enumerate}[nosep]
    \item Sensors measure voltage, current, temperature, and irradiance.
    \item ESP32 reads, validates, timestamps, and publishes data via MQTT.
    \item Backend subscribes to MQTT, validates payloads (Pydantic), and stores in the database.
    \item Preprocessing pipeline cleans data and engineers features.
    \item ML models produce consumption and PV forecasts.
    \item Optimization engine generates the next-day battery and load schedule.
    \item Dashboard displays real-time status, forecasts, schedules, and alerts.
\end{enumerate}

The simulation module can replace the ESP32 at the MQTT level, publishing synthetic data in the same JSON format. This makes the entire software stack hardware-agnostic.

\section{Technology Stack}

Table~\ref{tab:stack} summarizes the technology choices and their justifications.

\begin{table}[H]
\centering
\caption{Technology stack summary.}
\label{tab:stack}
\begin{tabularx}{\textwidth}{llX}
\toprule
\textbf{Component} & \textbf{Technology} & \textbf{Justification} \\
\midrule
Microcontroller & ESP32 DevKit V1 & Built-in Wi-Fi, dual-core, PlatformIO support \\
Communication & MQTT v3.1.1 & Lightweight pub/sub, IoT standard \\
Backend & Python 3.13 + FastAPI & Async, Pydantic validation, auto-docs \\
Database & SQLite / PostgreSQL & Zero-config dev; scalable production \\
ML & scikit-learn, XGBoost & Proven tabular ML; no unnecessary DL \\
Optimization & PuLP + CBC & Free MILP solver, sufficient for 96-step horizon \\
Dashboard & Streamlit & Native Python, rapid prototyping \\
\bottomrule
\end{tabularx}
\end{table}
